% =====================================================================
% Section 2: Related Work
% =====================================================================
%
% Status: Draft §2 Related Work v5 (2026-05-17)
%   v1-v3: see archive/
%   v4: GPT review #3 — added 100w on corpus-wide vs query-local
%       boundary (§2.1) and on avoiding separately trained reranker
%       (§2.3); 0 prose colons/semicolons.
%   v5 (this version): borrow NeocorRAG-style structural framing —
%        - tighten §2.2 to name "evidence chains" terminology used by
%          NeocorRAG, aligning the contrast surface with their paper.
%        - add a closing Summary paragraph that ties the three routes
%          back to the §1 conversion-problem framing (connected and
%          covering evidence list), mirroring NeocorRAG §2 tie-back.
%        - keep three-paragraph backbone untouched.
%
% Structure: three-paragraph form + closing summary
%   §2.1 Graph-based Retrieval for Multi-hop QA
%   §2.2 KG Reasoning and Query-time Chain Mining
%   §2.3 Personalized PageRank and Coverage-aware Reranking
%   §2.4 Summary (tie-back to §1 framing)
%
% =====================================================================

\section{Related Work}
\label{sec:related}

\paragraph{Graph-based retrieval for multi-hop QA.}
Standard retrieval-augmented generation methods rank passages by dense
or lexical similarity to the question
\citep{karpukhin2020dpr,robertson2009bm25}, which is effective for
single-hop questions but tends to miss bridge passages whose surface
similarity to the question is low. Recent work augments this pipeline
with explicit graph structure to improve multi-hop coverage.
HippoRAG \citep{jimenez2024hipporag} and HippoRAG 2
\citep{gutierrez2025hipporag} extract OpenIE facts from the corpus and
run personalized PageRank over a corpus-wide graph of passage and
phrase nodes seeded from the question. PropRAG \citep{wang2025proprag}
replaces triples with propositions and discovers multi-hop evidence by
beam search over a proposition graph, motivated by the observation that
triples can lose context. GraphRAG \citep{edge2024graphrag} builds
community-summary graphs over the corpus and supports retrieval through
hierarchical summary aggregation, while hypergraph-based methods such
as HGRAG \citep{hgrag2024} model entities and passages within a
hypergraph and propagate relevance across entity-level and
passage-level representations. Other recent variants explore
hierarchical multi-granular graphs \citep{hu2026mgrangrag},
sentence-level graph structures \citep{liang2026sentgraph}, and
linear CPU-only PageRank-style retrieval \citep{wang2026sprig}.
These methods mainly differ in the graph representation and in the
search operator used over the pre-built structure, ranging from
corpus-wide PageRank to beam search over propositions and hierarchical
summary aggregation. Many methods in this line rely on a pre-built
global structure or a summary hierarchy, with question conditioning
entering through seeds, retrieval scores, or a subsequent search
procedure. \methodname{} builds on the same OpenIE-based indexing line
\citep{angeli2015leveraging,gutierrez2025hipporag}, but preserves
document-grounded fact units for query-local propagation and then
applies relation-grounded expansion and coverage-aware reranking on top
of the local ranking. The activated region itself is a function of the
question, which makes the propagation behavior dependent on what the
question asks rather than only on which seeds are scored highly.

\paragraph{KG reasoning and query-time chain mining.}
A second line of work uses extracted relations as a reasoning substrate
that is traversed or updated at query time. ToG \citep{sun2024tog}
performs LLM-guided traversal over a pre-built knowledge graph,
selecting paths step by step, while related KG-RAG approaches
\citep{baek2023kaping} retrieve subgraphs and pass them to a language
model as context. Relink \citep{huang2026relink} addresses KG
incompleteness by dynamically reconstructing query-specific evidence
graphs and repairing missing facts as needed. NeocorRAG
\citep{peng2026neocorrag} keeps documents as the reader unit but
mines \emph{evidence chains} from retrieved text through an
activation-based candidate search followed by constrained decoding,
using the resulting chains to filter and re-present passages to the
reader. Graph-hop retrieval methods \citep{zhu2024reasongraphqa}
construct hop-wise evidence over textual databases. Interleaved
chain-of-thought retrieval methods such as IRCoT
\citep{trivedi2023ircot} and Self-Ask
\citep{press2023selfask} similarly drive retrieval through
language-model-generated reasoning traces, alternating between
generation and lookup. In contrast, \methodname{} keeps the OpenIE
structure on the indexing side. The offline graph is neither traversed
by an LLM at query time nor regenerated per question, and chain-aware
retrieval is computed by graph propagation rather than by query-time
chain generation over retrieved text. The reader consumes
natural-language passages, and any query-time language model use is
restricted to extracting question-side anchors and requirements.

\paragraph{Personalized PageRank and coverage-aware reranking.}
\methodname{}'s retrieval and enhancement components draw on
long-standing techniques in information retrieval. Personalized
PageRank \citep{haveliwala2002topic,jeh2003scaling} provides a
principled way to bias graph propagation toward a query, and has been
used in graph-augmented retrieval to score nodes given a seed
distribution. The propagation view is naturally complemented by a
selection view. Even when graph propagation has found a strong
candidate set, the order in which candidates enter the final list
determines what the reader actually sees. MMR \citep{carbonell1998mmr}
introduced diversity-aware reordering through marginal relevance
against an already-selected set, and submodular set functions
\citep{lin2011submodular} formalize coverage and diminishing returns
when selecting summary or evidence sets. Both perspectives motivate the
saturation behavior we adopt for $\mathrm{cov}_q$. Pointwise and
pairwise neural rerankers provide strong passage-level relevance
estimates, but they usually score candidates independently from the
current evidence set, which makes requirement-level coverage less
explicit. \methodname{} adapts both the graph-propagation and the
coverage-aware perspectives to the multi-hop retrieval setting.
Propagation is restricted to a query-induced subgraph and runs over
OpenIE fact units rather than over passage or word nodes directly,
while the coverage signal is computed under a noisy-OR model over
question-side entity-relation requirements rather than over query
tokens or word-level features. The two enhancement steps reuse the
same offline graph and the same dense embeddings as the local
propagation, so retrieval and reranking are aligned to the question's
relational structure rather than to surface similarity alone. This
design also avoids introducing a separately trained reranker on top of
the retrieval graph, since the saturation behavior of the noisy-OR
coverage and the question-conditioned propagation already provide
complementary signals at evidence-set granularity.

\paragraph{Summary.}
Existing GraphRAG systems improve cross-document retrieval through
richer graph representations, query-time path construction, or
chain-based filtering of retrieved text, and KG-RAG and
chain-of-thought retrieval methods further sharpen reasoning
trajectories. These approaches mostly target either how to build a
better corpus-level structure or how to surface explicit reasoning
chains from retrieved content, while the question of how to convert
graph-level relevance into a short reader-facing list that is both
connected across documents and covering of question-side requirements
receives less direct attention. \methodname{} addresses this
conversion problem within the GraphRAG family. It combines query-local
propagation over an induced fact-and-passage graph with
relation-aware evidence admission and coverage-aware reranking on top
of an offline OpenIE graph.
